\documentclass[12pt]{article}
\usepackage[utf8]{inputenc}
\usepackage[T1]{fontenc}
\usepackage{hyperref}
\title{Wikalkan Language Resources for K-12 Curriculum Development}
\author{Compiled by Claude Sonnet 5.5 (claude-sonnet-5-5) for Andr\'as Kornai}
\date{}
\begin{document}
\maketitle

\section{Introduction}
Wikalkan (Wik-Alken, Wik-Ngatharr, Wik Elken; ISO 639-3 wik) is a language of the Wik peoples of Cape York Peninsula, Queensland. Sources are inconsistent about how the name maps to the code and to the neighbouring Wik-Ngathan (wig). One page cites 86 speakers in 1981 for Wik-Ngatharr. Another reports only 3 speakers (2016 census) for Wik-Ngathan. Both figures should be treated cautiously.

\section{Primary Educational Resources}
No primer, textbook, or school programme was found for Wikalkan.

\section{Linguistic and Reference Materials}
Wikipedia pages: \url{https://en.wikipedia.org/wiki/Wik_Elken}, \url{https://en.wikipedia.org/wiki/Wik_peoples}, \url{https://en.wikipedia.org/wiki/Wik-Ngathan_language}, and the ISO 639:wik page \url{https://wikipedia.com/wiki/ISO_639:wik}. No dictionary or grammar specific to Wikalkan was found.

\section{Religious and Community Resources}
The search found no scripture or community organisation specific to Wikalkan. A general overview of the neighbouring Wik Mungkan is at \url{https://www.everyculture.com/Oceania/Wik-Mungkan-Orientation.html}.

\section{Digital and Audio Resources}
No audio, video, or corpus was found.

\section{Curriculum Development Context}
Documentation found is thin and mostly encyclopedic. Before development begins, the language's identity and community custodians should be confirmed, for example through AIATSIS AUSTLANG and community consultation. The search did not reach those sources.
\end{document}